\section{A residue calculation of every stability weight}\label{sec:indices}

The non-simple infinity index, rather than its multiplier or fixed
multiplicity, controls the critical weighted sum. Let $F$ be a rational map
fixing infinity and consider $\eta=dz/(z-F(z))$. A finite pole of $F$ is a
zero of this form. Its other possible finite poles are fixed points, and at a
simple fixed point its residue is $1/(1-F'(z))$. For
$g(w)=1/F(1/w)$ a direct coordinate transformation gives
\begin{equation}\label{eq:form-change}
 \eta=\left(\frac1{w-g(w)}-\frac1w\right)dw.
\end{equation}
Let $I_\infty(F)$ be the coefficient of $w^{-1}$ in $1/(w-g(w))$. The
residue theorem on the sphere gives
\begin{equation}\label{eq:index}
 \sum_{\substack{F(z)=z\\z\text{ finite}}}
 \Res_z\frac{d\zeta}{\zeta-F(\zeta)}+I_\infty(F)=1.
\end{equation}
The subtractive $-1/w$ in~\eqref{eq:form-change} is necessary; it is the source
of the constant on the right.

Apply this to $F=T^n$. All finite fixed points are simple by
Section~\ref{sec:dynamics}, including the subcritical off-real pair. For
$a\ne1$, the index is $I_\infty(T^n)=1/(1-a^{-n})$. At $a=1$,
\eqref{eq:jet} instead gives
\begin{equation}\label{eq:critical-index}
 \frac1{w-g_1^n(w)}=-\frac1{nbw^3}
   +\frac{3n-1}{2n}\frac1w+O(w),\qquad
 I_\infty(T^n)=\frac{3n-1}{2n}.
\end{equation}
In particular, the fixed multiplicity $3$ is not this index.

Changing signs in the finite terms of~\eqref{eq:index} yields
\begin{equation}\label{eq:finite-sum}
 \sum_{\substack{T^n(z)=z\\z\text{ finite}}}
 \frac1{(T^n)'(z)-1}=I_\infty(T^n)-1.
\end{equation}
Away from criticality the right side is $1/(a^n-1)$. For $a<1$, the two
nonreal fixed points have multiplier $q^n$ and each contributes
$1/(q^n-1)$, so they must be subtracted. For $a\ge1$ all finite fixed points
are real. Substitution proves~\eqref{eq:tau-main}, including the critical value
$(n-1)/(2n)$.

All weights are nonnegative, as also follows from their original positive
denominators. For $a\le1$, $N_1=\tau_1=0$. At criticality the two-cycle
$\{\sqrt{b/2},-\sqrt{b/2}\}$ has multiplier $9$, so its two points contribute
$2/(9-1)=1/4$, consistent with the all-iterate formula at $n=2$. This example
is an exact consequence, not a substitute for the index argument.
